% !TEX root = ./../main.tex
\section{Introduction}
\label{sec:introduction}
Diffusion probabilistic models have achieved remarkable success in generating high-quality images~\citep{ho2020denoising,song2021scorebased,dhariwal2021diffusion,rombach2022high}.
On top of the success in the image domain, there has been rapid progress in the generation of videos~\citep{ho2022video,singer2022make,zhou2022magic,wang2023modelscope}.
%
Despite the progress, long video generation still lags behind compared to image generation.
One reason is that video diffusion models~(VDMs) often consider a video as a single 4D tensor with an additional axis corresponding to time, which prevents the models from generating videos at scale.
An intuitive approach to generating a long video is autoregressive generation, which iteratively predicts a future frame given the previous ones.
However, in contrast to the transformer-based models~\citep{hong2023cogvideo,villegas2023phenaki}, diffusion-based models cannot directly adopt the autoregressive generation strategy due to the heavy computational costs incurred by iterative denoising steps for a single frame generation.
Instead, several recent works~\citep{ho2022video,he2022latent,voleti2022mcvd,luo2023videofusion,chen2023seine,blattmann2023align} adopt a chunked autoregressive generation strategy, which predicts several frames in parallel conditioned on few preceding ones, consequently reducing computational burden.
While these approaches are computationally tractable, they often leads to temporal inconsistency and discontinuous motion, especially between the chunks predicted separately, because the model captures a limited temporal context available in the last few---only one or two in practice---frames.

To address the limitation, we propose a novel inference technique, FIFO-Diffusion, which realizes arbitrarily long video generation without training based on a pretrained video generation model for short clips.
Our approach effectively alleviates the limitations of the chunked autoregressive method by enabling every frame to refer to a sufficient number of preceding frames. 
Our approach generates frames through diagonal denoising~(\cref{subsec:diagonal}) in a first-in-first-out manner using a queue, which contains a sequence of frames with different---monotonically increasing---noise levels over time.
At each step, a completely denoised frame at the head is popped out from the queue while a new random noise image is pushed back at the tail.
Diagonal denoising offers both advantage and disadvantage; noisier frames benefit from referring to cleaner ones while the model may suffer from training-inference gap because video models are generally trained to denoise frames with the same noise level.
To overcome this trade-off and embrace the advantage of diagonal denoising, we propose latent partitioning~(\cref{subsec:latent_partitioning}) and lookahead denoising~(\cref{subsec:look_ahead_denoising}). 
Latent partitioning reduces training-inference gap by narrowing the range of noise levels in to-be-denoised frames and enables inference with finer steps.
Lookahead denoising allows to-be-denoised frames to reference cleaner frames, thereby performing more accurate noise prediction.
Furthermore, both latent partitioning and lookahead denoising offer parallelizability on multiple GPUs.


\vspace{2mm}
Our main contributions are summarized below.
\begin{itemize}
    \item We propose FIFO-Diffusion through diagonal denoising, which is a training-free video generation technique for VDMs pretrained on short clips.
    Our approach denoises images with different noise levels for seamless video generation, enabling us to generate arbitrarily long videos.
    \item We introduce latent partitioning and lookahead denoising, which respectively reduce the training-inference gap incurred by diagonal denoising and allow the reference to less noisy frames for denoising, improving generation quality.
    \item FIFO-Diffusion requires a constant amount of memory regardless of the length of the generated videos given a baseline model. It is straightforward to run FIFO-Diffusion in parallel on multiple GPUs.
    \item Our experiments on four strong baselines, based on the U-Net~\citep{ronneberger2015unet} or DiT~\citep{peebles2023dit} architectures, show that FIFO-Diffusion generates extremely long videos including natural motion without degradation on quality over time.
\end{itemize}